
Lista curada y \textbf{no exhaustiva}. Prioriza marcos conceptuales y buenas prácticas; actualice enlaces por edición.

\subsection*{IA agéntica y sistemas de agentes}
\begin{itemize}
\item Documentación conceptual de agentes / tool use de proveedores mayores --- comparar definiciones de \emph{tools}, \emph{handoffs} y memoria.
\item Artículos introductorios sobre \emph{human-in-the-loop} y orquestación de workflows.
\item Guías de \emph{prompting} estructurado y evaluación de salidas.
\end{itemize}

\subsection*{Diseño instruccional y contenido académico}
\begin{itemize}
\item Materiales institucionales del Tec sobre diseño de cursos *(usar los vigentes del campus / CEDDIE)*.
\item Referencias de alineación constructiva (objetivos--actividades--evaluación).
\item Guías de accesibilidad de contenidos digitales (WCAG resumida para docentes).
\end{itemize}

\subsection*{Integridad académica y uso responsable de IA}
\begin{itemize}
\item Políticas del Tecnológico de Monterrey sobre integridad académica y uso de IA generativa --- \ceddie.
\item Declaraciones de editores sobre divulgación de uso de IA en manuscritos.
\item Lecturas sobre alucinaciones, citas fabricadas y verificación de fuentes.
\end{itemize}

\subsection*{Escritura académica y revisión}
Manuales de estilo relevantes (APA, IEEE, Chicago, etc.); guías breves de revisión por pares.

\subsection*{Cómo usar esta lista}
\begin{tabularx}{\textwidth}{@{}lX@{}}
\toprule
Momento & Uso sugerido \\
\midrule
Sesión 1 & 1 lectura corta de integridad + 1 overview de agentes \\
Sesiones 2--3 & Consulta asíncrona según el artefacto \\
Sesión 7--8 & Política de divulgación de IA en publicación \\
\bottomrule
\end{tabularx}
